% A self-contained copy of the proof of Proposition prop:kappa0 and of the m = 0 correction,
% extracted from paper/level-prime-kappa.tex (section "Where the correction comes from") so that it
% can be read and commented on in Overleaf independently of the rest of the preprint.
%
% NOTHING HERE IS AGREED YET. The argument has been reviewed by agents on both sides and corrected
% twice; it still needs Eran's and Sachi's agreement before the code or the tests describe
% Proposition prop:kappa0 as proved.
%
% Numbering is local to this file, so references will not match the preprint's.
\documentclass[11pt,reqno]{amsart}
\usepackage[margin=1.15in]{geometry}
\usepackage{amsmath,amssymb,amsthm,mathtools}
\usepackage{booktabs}
\usepackage[colorlinks=true,linkcolor=blue,citecolor=blue,urlcolor=blue]{hyperref}
\theoremstyle{plain}
\newtheorem{theorem}{Theorem}
\newtheorem{proposition}[theorem]{Proposition}
\newtheorem{lemma}[theorem]{Lemma}
\newtheorem{corollary}[theorem]{Corollary}
\theoremstyle{definition}
\newtheorem{observation}[theorem]{Observation}
\newtheorem{remark}[theorem]{Remark}
\DeclareMathOperator{\Diff}{Diff}
\DeclareMathOperator{\ord}{ord}
\DeclareMathOperator{\disc}{disc}
\DeclareMathOperator{\tr}{tr}
\DeclareMathOperator{\nrd}{nrd}
\newcommand{\Q}{\mathbb{Q}}
\newcommand{\Z}{\mathbb{Z}}
\newcommand{\C}{\mathbb{C}}
\newcommand{\Ldual}{L^{\vee}}
\newcommand{\Lm}{L_{-}}
\newcommand{\Lp}{L_{+}}
% references to the parent preprint, so this file compiles alone
\newcommand{\parentref}[1]{\textup{[#1 in the preprint]}}
\begin{document}
\title{The dropped $m=0$ coefficient at a split level prime}
\author{}
\date{\today}
\maketitle

\noindent\textbf{Status.} This is the argument for Proposition~\ref{prop:kappa0} below and for the
$m=0$ correction it implies, extracted so it can be read on its own. It has been reviewed twice by
agents and corrected both times; it is \emph{not} yet agreed, and nothing in the accompanying code
or tests describes it as proved. Objects and notation ($\Lm$, $\Lp$, $\lambda$, the Eisenstein
coefficients $A_\mu(s,m,v)$, the multiplier of Observation~obs:mult) are those of the preprint's
Sections~sec:chain and sec:data; cross-references into the preprint are left as plain labels.

\subsection*{The literature's zero, and where it comes from}

In the published form of Schofer's theorem \cite[Thm.~1.1]{Schofer} the sum runs over all cosets
$\eta$ and all $m\ge0$; the vanishing of the $m=0$ terms at $\eta\neq0$ is
\cite[Lemma~2.21]{Schofer}, ``for $\mu\neq0$, $b_\mu(0,t)=0$'', stated there for a general lattice.
Its proof reads: ``Theorem~5.2 of [17] implies $W_{0,p}(s,\varphi_{\mu,p})=0$ if $\varphi_{\mu,p}$
is not the characteristic function of the local lattice'', where [17] is an unpublished preprint
(T.~Yang, 2001), so the step cannot be checked as stated. The nearest published statement we can
check is \cite[Prop.~5.2]{KY}, and it is proved for the ring of integers of a quadratic field,
where every nonzero coset of the relevant lattice is anisotropic and the conclusion is immediate.
At an Eichler order of level $N>1$ the lattice $\Lm$ is $N$-scaled at the level prime and has
nonzero \emph{isotropic} cosets, and there Lemma~\ref{lem:WN} computes the local factor to be a
nonzero constant. So what we are contradicting is not a verifiable proof but a lemma whose only
published analogue is proved under a hypothesis these lattices violate; Yang's \cite[(11)]{Yang}
inherits the conclusion.
Everything else in the chain --- the constant-term factorisation \cite[Thm.~2.4]{KY}, the vanishing
of the incoherent series at $s=0$, and the local structure of $\Lm$ \cite[Lemma~18]{Yang} --- is
used unchanged. One point of bookkeeping decides the shape of the final statement: the quantity
set to zero enters the Stokes argument \cite[(33)]{Schofer} additively through $b_\mu(0,t)$, which
is constant in $t$ here, so replacing it by $\lim_tb_\mu(0,t)$ is legitimate; but it must be
replaced \emph{wherever} Yang's (12) invokes it, and (12) calls $\kappa^-_\nu(0)$ not only at
$m=0$ but at every $m>0$ through the vectors $x$ on the CM line with $Q(x)=m$.
Proposition~\ref{prop:mult} keeps all of them.

Throughout, $N$ is prime, $d<0$ is a fundamental discriminant with $N\nmid d$ (so $N$ splits in
$k=\Q(\sqrt d)$ by Eichler's condition, the hypothesis $\mathrm{CM}(d)\neq\varnothing$ of
\cite[Thm.~B]{GY}), $\lambda=\varphi(\sqrt d)$, $\langle\alpha,\beta\rangle=\tr(\alpha\beta')$ so
that $\langle\lambda,\lambda\rangle=2\lvert d\rvert$, $Q=\langle\cdot,\cdot\rangle/2$,
$\varphi_\nu=\mathrm{char}(\nu+\widehat{\Lm})$, and $A_\nu(s,0,v)=b_\nu(0,v)\,s+O(s^2)$.

\begin{lemma}[the lattice at the level prime]\label{lem:LN}
$L_{-,N}$ is the $N$-scaled hyperbolic plane, $N\begin{pmatrix}0&1\\1&0\end{pmatrix}$ up to a
unit; $L_{-,N}^\vee/L_{-,N}\cong(\Z/N)^2$ with $Q(x,y)=xy/N\bmod\Z$, and the nonzero isotropic
cosets are the $2(N-1)$ nonzero vectors on the two lines $x=0$, $y=0$. Moreover
$L_N=L_{+,N}\oplus L_{-,N}$.
\end{lemma}
\begin{proof}
\cite[Lemma~18(1)]{Yang}: for odd $N\nmid d$ the Gram matrix of $L_{-,N}$ is
$N\,\mathrm{diag}(\epsilon_1,\epsilon_2)$ with $\epsilon_1\epsilon_2=-d$, and $\chi_d(N)=1$ makes
$-\epsilon_1\epsilon_2$ a square, so the plane is split. For $N=2$ (then $d\equiv1\bmod8$)
\cite[Lemma~18(3)]{Yang} gives $2\epsilon\begin{pmatrix}0&1\\1&0\end{pmatrix}$ directly (Yang
writes out the $2$-adic proof only for $2\nmid DN$; we checked this case: $L_{-,2}$ has basis
$\begin{pmatrix}-c_3/c_1&1\\0&c_3/c_1\end{pmatrix}$,
$\begin{pmatrix}-c_2/c_1&0\\2&c_2/c_1\end{pmatrix}$ in the notation below, with Gram matrix twice
an even form with odd off-diagonal entry and determinant $\equiv-1\bmod8$).

For the direct sum it suffices that $\langle\lambda,L_N\rangle\subset\langle\lambda,\lambda\rangle
\Z_N$: then the orthogonal projection $x\mapsto\frac{\langle x,\lambda\rangle}{\langle\lambda,
\lambda\rangle}\lambda$ maps $L_N$ into $\Z_N\lambda=L_{+,N}$ and $x-\mathrm{pr}(x)\in L_N\cap
\lambda^\perp=L_{-,N}$. For odd $N$, $2\lvert d\rvert$ is an $N$-unit. For $N=2$:
$\mathcal O\otimes\Z_2=\begin{pmatrix}\Z_2&\Z_2\\2\Z_2&\Z_2\end{pmatrix}$ (remark in the proof of
\cite[Lemma~18]{Yang}), $\lambda=\begin{pmatrix}c_1&c_2\\2c_3&-c_1\end{pmatrix}$, and optimality
of the embedding of $\mathcal O_d$, $d\equiv1\bmod8$, means $(1+\lambda)/2\in\mathcal O\otimes
\Z_2$, i.e.\ $c_1$ odd and $c_2,c_3$ even; for $\alpha=\begin{pmatrix}a&b\\2c&-a\end{pmatrix}\in
L_2$, $\langle\alpha,\lambda\rangle=-\tr(\alpha\lambda)=-2(ac_1+bc_3+cc_2)\in2\Z_2=
2\lvert d\rvert\Z_2$.
\end{proof}

Note that $L_{+,N}=\Z_N\lambda$ is unimodular for odd $N$ but not for $N=2$, where
$L_{+,2}^\vee/L_{+,2}\cong\Z/2$ is generated by $\lambda/2$ with $Q(\lambda/2)=\lvert d\rvert/4
\notin\Z_2$.

\begin{lemma}[the $m=0$ local factor at $N$]\label{lem:WN}
With $X=N^{-s}$ and $\gamma_N$ the local splitting index of \cite[(4.3)]{KY}, for $\nu\in
L_{-,N}^\vee/L_{-,N}$,
\[
  W_{0,N}(s,\varphi_0)=\frac{\gamma_N}{N}\cdot\frac{1+(N-2)X}{1-X},\qquad
  W_{0,N}(s,\varphi_\nu)=\frac{\gamma_N}{N}\ \text{($\nu\neq0$ isotropic)},\qquad
  W_{0,N}(s,\varphi_\nu)=0\ \text{($\nu\neq0$ anisotropic)}.
\]
\end{lemma}
\begin{proof}
The first is \cite[(15)]{Yang} with $\chi_d(N)=1$ for odd $N$, where $\frac{1-X/N}{N-X}=\frac1N$,
and \cite[(16)]{Yang} for $N=2$. For the others apply \cite[(4.4) and Thm.~4.3]{KY} ($N$ odd; KY's
matrix $S$ is the matrix of the bilinear form with $Q=S/2$, as is Yang's Gram matrix, so
$S=N\,\mathrm{diag}(\epsilon_1,\epsilon_2)$, $l_1=l_2=1$, $s_0=n/2-1=0$) to $\nu=(a/N,b/N)$ with
$m=0$: $Q(\nu)=(\epsilon_1a^2+\epsilon_2b^2)/(2N)$ lies in $\Z_N$ exactly when $\nu$ is isotropic
(for odd $N$ the $2$ is a unit and may be dropped), so
$W=0$ otherwise by Thm.~4.3(i); when it is, $-\epsilon_1\epsilon_2$ being a square forces $a,b$
both units, so $H_\nu=\varnothing$, $K_0(\nu)=\min(l_i+\ord_N\nu_i)=0$, and $a_\nu=\ord_N(t_\nu)
\ge0=K_0$ puts us in the second case of Thm.~4.3(ii), whose sum over $0<k\le K_0$ is empty:
$W_N(s,0,\nu)=1$. For $N=2$, \cite[Thm.~4.4]{KY} with the Gram matrix $2\epsilon\begin{pmatrix}
0&1\\1&0\end{pmatrix}$, i.e.\ KY's (4.9) with $H=N=0$, $M=1$, $m_1=1$: $\nu'=(a/2,b/2)\notin
\Z_2^2$ gives $M_\nu=\varnothing$, $K_0=m_1+\ord_2\nu'=0$, both sums in (ii) are empty, and
$W_2(s,0,\nu)=1$ when $Q(\nu)=ab/2\in\Z_2$, i.e.\ at the two nonzero isotropic cosets. In every
case the normalising constant of \cite[(4.4)]{KY} is $\gamma_N\lvert\det S\rvert_N^{1/2}=
\gamma_N/N$, the same for all $\nu$, so it cancels in the ratio used below. \textup{Lemma~lem:m0loc of the preprint}
is the same object in a counting normalisation (its unimodular calibration is Thm.~4.3 with
$l_i=0$), and its third and fourth rows are these values.
\end{proof}

\begin{lemma}[no other nonzero coset contributes]\label{lem:Wp}
For $p\neq N$ every nonzero $\nu_p\in L_{-,p}^\vee/L_{-,p}$ is anisotropic, $Q(\nu_p)\notin\Z_p$.
Consequently $W_{0,p}(s,\varphi_{\nu,p})=0$ for such $\nu_p$, and a nonzero isotropic coset of
$\Lm^\vee/\Lm$ is supported at $N$.
\end{lemma}
\begin{proof}
$W_{0,p}=0$ follows from anisotropy by \cite[Thm.~4.3(i)]{KY} ($p$ odd) and \cite[Thm.~4.4(i)]{KY}
($p=2$). Anisotropy is read off \cite[Lemma~18]{Yang}; $d$ fundamental gives $p\,\|\,d$ for odd
$p\mid d$.

$p$ odd, $p\nmid dDN$: unimodular. $p$ odd, $p\mid d$: Gram $\mathrm{diag}(\epsilon_1,\epsilon_2)$
with, say, $\epsilon_2=p\epsilon_2'$; $L^\vee/L\cong\Z/p$ generated by $(0,1/p)$, and
$Q(0,b/p)=\epsilon_2'b^2/(2p)\notin\Z_p$ for $p\nmid b$. $p$ odd, $p\mid D$, $p\nmid d$: Gram
$p\,\mathrm{diag}(\epsilon_1,\epsilon_2)$ with $\epsilon_1\epsilon_2=-d$; $p$ is inert in $k$, so
$d$ is a non-square mod $p$ (the Hilbert-symbol condition $(-\mu_1,-\mu_2p)_p=-1$ in Yang's
proof), and $Q(a/p,b/p)=(\epsilon_1a^2+\epsilon_2b^2)/(2p)\in\Z_p$ forces $a\equiv b\equiv0$.

$p=2\nmid d$. If $d\equiv1\bmod8$ then $2\nmid D$ and, as $2\neq N$, Lemma~18(3) gives
$\epsilon\begin{pmatrix}0&1\\1&0\end{pmatrix}$, unimodular. If $d\equiv5\bmod8$, Lemma~18(4) gives
$\epsilon\begin{pmatrix}2&1\\1&2\end{pmatrix}$ for $2\nmid D$, unimodular, and $2\epsilon
\begin{pmatrix}2&1\\1&2\end{pmatrix}$ for $2\mid D$, where $L^\vee/L\cong(\Z/2)^2$ and
$Q(a/2,b/2)=\epsilon(a^2+ab+b^2)/2$ with odd numerator for $(a,b)\not\equiv(0,0)$. (This case is
again one Yang only sketches; with $(1+\lambda)/2$ in the maximal order, $a,b,c$ are all odd in
$\lambda=aI+bJ+cIJ$ below, and the elimination below gives diagonal entries $\equiv2\bmod8$, a
unit off-diagonal and determinant $\equiv3\bmod8$.)

$p=2\mid d$: Lemma~18(2) gives $2\,\mathrm{diag}(\epsilon_1,\epsilon_2)$ with
$\epsilon_1\epsilon_2=-d/4$, $Q(x,y)=\epsilon_1x^2+\epsilon_2y^2$, and $d$ fundamental means
$d/4\equiv2,3\bmod4$. If $d\equiv12\bmod16$ both $\epsilon_i$ are units with $\epsilon_1\equiv
\epsilon_2\bmod4$, and the three nonzero cosets of $(\Z/2)^2$ have $Q=\epsilon_1/4$,
$\epsilon_2/4$, $(\epsilon_1+\epsilon_2)/4\equiv\epsilon_1/2$. If $d\equiv8\bmod16$, say
$\epsilon_2=2\epsilon'$; $L^\vee/L\cong\Z/2\times\Z/4$ and $Q(a/2,b/4)=\epsilon_1a^2/4+
\epsilon'b^2/8$ has $2$-adic valuation $-3$ for $b$ odd, $-2$ or $-1$ for $b\equiv2\bmod4$, and
equals $\epsilon_1/4$ for $b\equiv0$, $a$ odd. None is in $\Z_2$.

Yang writes out the proof of Lemma~18(2) only for $2\nmid DN$; the case $2\mid N$, $2\mid d$ does
not occur here, and the case $2\mid D$, $2\mid d$ is proved as follows. $B\otimes\Q_2=
\bigl(\frac{-1,-1}{\Q_2}\bigr)$ with maximal order $\Z_2+\Z_2I+\Z_2J+\Z_2(1+I+J+IJ)/2$, whose
trace-zero part is $L_2=\Z_2I+\Z_2J+\Z_2IJ$ with $Q=x^2+y^2+z^2$. Since $d\equiv0\bmod4$,
$\mathcal O_d=\Z[\sqrt d/2]$ and $\Lp=\Z\lambda_0$ with $\lambda_0=\lambda/2=aI+bJ+cIJ$,
$a^2+b^2+c^2=m:=-d/4\equiv1$ or $2\bmod4$, and $L_{-,2}=\{xI+yJ+zIJ:ax+by+cz=0\}$. If
$m\equiv1\bmod4$ exactly one of $a,b,c$ is odd, say $a$; eliminating $x=-(by+cz)/a$,
$Q|_{L_{-,2}}=u_1y^2+2w\,yz+u_2z^2$ with $u_1=b^2/a^2+1$, $u_2=c^2/a^2+1$ units and $w=bc/a^2\in
4\Z_2$; $y'=y+(w/u_1)z$ is an integral change of basis and $Q=u_1y'^2+(u_2-w^2/u_1)z^2$ with
$u_1u_2-w^2=m/a^2$. If $m\equiv2\bmod4$ two of $a,b,c$ are odd, say $a,b$; then $u_1\in2\Z_2^\times$,
$u_2$ is a unit, $w\in2\Z_2$, and completing the square on $z$ instead gives $Q=u_2z'^2+
(u_1-w^2/u_2)y^2$, product again $m/a^2$. Rescaling one basis vector by the unit $a$ makes the
product $-d/4$ exactly, which is Lemma~18(2).
\end{proof}

\begin{proposition}[the dropped coefficient]\label{prop:kappa0}
Let $N$ be prime, $N\nmid d$, and $\nu\in\Lm^\vee/\Lm$ a nonzero isotropic coset. Then
$b_\nu(0,v)=-\log N/(N-1)$ for every $v$, hence
\[
  \kappa^-_\nu(0)\;:=\;\lim_{v\to\infty}b_\nu(0,v)\;=\;-\frac{\log N}{N-1}.
\]
\end{proposition}
\begin{proof}
By \cite[Thm.~2.4]{KY}, as used in \cite[proof of Lemma~2.21]{Schofer} and \cite[(14)]{Yang}, the
constant term is $E_0(\tau,s;\varphi_\nu,+1)=v^{s/2}\varphi_\nu(0)+W_{0,\infty}(\tau,s)\prod_p
W_{0,p}(s,\varphi_{\nu,p})$ with $W_{0,\infty}$ finite at $s=0$ (\cite[Prop.~2.3(iii)]{KY} with $l=1$; only finiteness is
used below). Put $\Phi(s)=E_0(\tau,s;\varphi_0,+1)-v^{s/2}$. For $\nu\neq0$, $\varphi_\nu(0)=0$, and
$\varphi_{\nu,p}=\varphi_{0,p}$ for $p\neq N$ (Lemma~\ref{lem:Wp}), so for $\mathrm{Re}\,s$ large,
hence by meromorphic continuation,
\[
  E_0(\tau,s;\varphi_\nu,+1)
   =\Phi(s)\,\frac{W_{0,N}(s,\varphi_\nu)}{W_{0,N}(s,\varphi_0)}
   =\Phi(s)\,\frac{1-N^{-s}}{1+(N-2)N^{-s}}
\]
by Lemma~\ref{lem:WN}. (The Euler product does not converge at $s=0$; $\Phi$ and the identity are
defined through Thm.~2.4 and continued. The pole of $W_{0,N}(s,\varphi_0)$ at $s=0$ is cancelled
by the missing Euler factor $1-N^{-s}$ of $L^S(s,\chi_d)$ there, and $W_{0,N}(s,\varphi_0)$ has no
zero near $s=0$.) The series $E(\tau,s;\varphi_0,+1)$ is holomorphic at $s=0$
(\cite[Cor.~2.5]{KY}) and incoherent, so it vanishes there (\cite{Yang}, the paragraph before
(11)); its constant term gives $\Phi(0)=-1$. The last factor vanishes at $s=0$ --- as it must,
$E(\tau,s;\varphi_\nu,+1)$ being incoherent --- with derivative $\log N/(N-1)$, so
$A_\nu(s,0,v)=-\frac{\log N}{N-1}s+O(s^2)$, with no $v$ remaining: the $v$-dependence of $\Phi$
enters at order $s$, multiplying a term already $O(s)$.

As a cross-check, \cite[(14)]{Yang} writes $\Phi(s)=-v^{-s/2}\frac{\Lambda(s,\chi_d)}
{\Lambda(s+1,\chi_d)}G(s)$ with $G=\prod_{p\mid r}G_p$ over the non-unimodular primes; the
functional equation of the completed $L$-function of the odd primitive character $\chi_d$ gives
$\Lambda(0)=\Lambda(1)$, and $G(0)=1$ is read off (15)--(17) factor by factor, so $\Phi(0)=-1$
again. The sign agrees with \cite[Lemma~20]{Yang}, where $-\frac{d}{ds}\log G_N\big|_{s=0}=\log N$
at the zero coset.
\end{proof}

\begin{lemma}[isotropic cosets of $\Ldual/L$]\label{lem:iso}
For $N$ prime the isotropic cosets of $\Ldual/L$ are supported at $N$ and number $2N-1$: under
$L_N=L_{+,N}\oplus L_{-,N}$ they are the cosets $0\oplus\nu$ with $\nu$ an isotropic coset of
$L_{-,N}^\vee/L_{-,N}$. All $2N-2$ nonzero ones carry the same constant term $c_\eta(0)$ of $F_f$.
\end{lemma}
\begin{proof}
A coset is isotropic iff each $p$-part is, and $\Ldual/L$ does not depend on $d$, so the
computation is on $L_p$ itself. For $p$ odd, $p\nmid DN$, $L_p$ is unimodular. For $p$ odd,
$p\mid DN$, the proof of \cite[Lemma~18]{Yang} gives a basis of $L_p$ with Gram matrix
$\mathrm{diag}(2\mu_1,2\mu_2p,2\mu_1\mu_2p)$, $\mu_i$ units, so $(\Ldual/L)_p\cong(\Z/p)^2$ with
$Q(ae_2/p+be_3/p)=\mu_2(a^2+\mu_1b^2)/p$: anisotropic iff $-\mu_1$ is a non-square mod $p$, i.e.\
iff $(-\mu_1,-\mu_2p)_p=-1$, which Yang's proof identifies with $p\mid D$; for $p\mid N$ it is the
hyperbolic plane over $\mathbb F_p$, consistent with Lemma~\ref{lem:LN}. For $p=2\nmid DN$, $L_2$
is the trace-zero part of $M_2(\Z_2)$, with Gram $(-2)\oplus\begin{pmatrix}0&-1\\-1&0\end{pmatrix}$
in the basis $\mathrm{diag}(1,-1),e_{12},e_{21}$, so $(\Ldual/L)_2\cong\Z/2$ generated by
$\tfrac12\mathrm{diag}(1,-1)$ with $Q=-1/4$. For $p=2\mid D$, $L_2=\Z_2I+\Z_2J+\Z_2IJ$ as in the
proof of Lemma~\ref{lem:Wp}, $(\Ldual/L)_2\cong(\Z/2)^3$ with $Q=(a^2+b^2+c^2)/4\in\{\tfrac14,
\tfrac12,\tfrac34\}$ on nonzero elements. At $p=N$ odd, $L_{+,N}$ is unimodular and
$(\Ldual/L)_N=L_{-,N}^\vee/L_{-,N}$, with $2N-1$ isotropic vectors. At $N=2$, the orthogonal direct
sum gives $(\Ldual/L)_2=L_{+,2}^\vee/L_{+,2}\oplus L_{-,2}^\vee/L_{-,2}\cong\Z/2\oplus(\Z/2)^2$; the
elements $\lambda/2\oplus\nu$ have $Q=\lvert d\rvert/4+Q(\nu)$ with $Q(\nu)\in\{0,\tfrac12\}\bmod
\Z$ and $\lvert d\rvert$ odd, never integral, so the isotropic elements are the $0\oplus\nu$ with
$\nu$ isotropic: again $2N-1=3$.

For the constant terms: an isometry $\sigma$ of the discriminant form permutes the $e_\eta$, fixes
$e_0$, and commutes with $\rho_L(S)$ and $\rho_L(T)$ because it preserves $Q$ and the pairing,
hence with every $\rho_L(\tilde\gamma)$; so $\sigma F_f=F_f$ and $c_{\sigma\eta}(0)=c_\eta(0)$. The
orthogonal group of the hyperbolic plane over $\mathbb F_N$ contains the torus $(x,y)\mapsto
(ax,a^{-1}y)$ and the swap, and acts transitively on the $2(N-1)$ nonzero isotropic vectors; these
isometries of $(\Ldual/L)_N$ extend by the identity on the other $p$-parts.
\end{proof}

\begin{proposition}[the restored terms of Theorem~B]\label{prop:mult}
Let $N$ be prime, $N\nmid d$, $d$ fundamental, and let $\nu$ run over the $2N-2$ nonzero isotropic
cosets of $\Lm^\vee/\Lm$, identified with those of $\Ldual/L$ by Lemma~\ref{lem:iso}. Then
$\kappa_\nu(0)=\kappa^-_\nu(0)=-\log N/(N-1)$, and the terms of \cite[Thm.~B]{GY} that
\cite[(11)]{Yang} sets to zero are, in total,
\[
  -\frac{\lvert\mathrm{CM}(d)\rvert}{4}\sum_{\nu}\kappa^-_\nu(0)
     \sum_{x\in\Ldual\cap\Q\lambda}c_{[x+\nu]}(-Q(x))
  \;=\;\frac{\lvert\mathrm{CM}(d)\rvert}{4}\,\frac{\log N}{N-1}
     \sum_{\nu}\ \sum_{x\in\Ldual\cap\Q\lambda}c_{[x+\nu]}(-Q(x)),
\]
with $[\,\cdot\,]$ the class in $\Ldual/L$. The $x=0$ term of each inner sum is $c_{[\nu]}(0)$, and
those $2N-2$ values agree (Lemma~\ref{lem:iso}), so the $x=0$ part of the total is
$\tfrac12\lvert\mathrm{CM}(d)\rvert c_\eta(0)\log N$; the $x\neq0$ terms are identified in
Lemma~\ref{lem:H}(3) and Remark~\ref{rem:xsum}, and are not assumed here to be independent of
$\nu$. In particular, if
\begin{itemize}
\item[(H)] $c_{[x+\nu]}(-Q(x))=0$ for every $x\neq0$ in $\Ldual\cap\Q\lambda$ and every nonzero
      isotropic $\nu$,
\end{itemize}
the correction is $\lvert\mathrm{CM}(d)\rvert\cdot\tfrac12c_\eta(0)\log N$, i.e.\
$\mathrm{mult}(f)\sum_{p\mid N/(N,d)}\log p$ per CM point with $\mathrm{mult}(f)=\tfrac12c_\eta(0)$,
which is \textup{Observation~obs:mult of the preprint}. (Theorem~B constrains the sum over $\mathrm{CM}(d)$; the
correction does not depend on the point, so dividing by $\lvert\mathrm{CM}(d)\rvert$ is the
convention of \textup{Section~sec:chain of the preprint}.)
\end{proposition}
\begin{proof}
Theorem~B sums $c_\eta(-m)\kappa_\eta(m)$ over $\eta\in\Ldual/L$ and $m\ge0$, with $\kappa_\eta(m)$
as in \textup{Section~sec:chain of the preprint}: $\kappa_\eta(m)=\sum_{\mu\in L/(\Lp+\Lm)}\sum_{x\in\eta_++\mu_++\Lp}
\kappa^-_{\eta_-+\mu_-}(m-Q(x))$, where $\eta=\eta_++\eta_-$ is the decomposition in
$\Lp^\vee\oplus\Lm^\vee\supset\Ldual$. The pairs $(\eta,\mu)$ run over $\Ldual/(\Lp\oplus\Lm)$
exactly once, and every factor depends only on $y=\eta+\mu$; the terms in which $\kappa^-$ is
evaluated at $0$ are those with $m=Q(x)$ ($\Lp$ is positive definite and $\kappa^-_\nu(m)=0$ for
$m<0$, so no other term involves $\kappa^-_\nu(0)$), and their total is
\[
  \sum_{y\in\Ldual/(\Lp\oplus\Lm)}\ \sum_{x\in y_++\Lp}c_{[y]}(-Q(x))\,\kappa^-_{[y_-]}(0)
  \;=\;\sum_{\nu\in\Lm^\vee/\Lm}\kappa^-_\nu(0)\sum_{\substack{x\in\Lp^\vee\\ x+\nu\in\Ldual}}
      c_{[x+\nu]}(-Q(x)),
\]
the inner sum over all $x$ in the coset, both signs included. By Lemma~\ref{lem:Wp} and
Proposition~\ref{prop:kappa0}, among nonzero $\nu$ the coefficient $\kappa^-_\nu(0)$ is nonzero
only for $\nu$ isotropic, supported at $N$, where it equals $-\log N/(N-1)$ (the $\nu=0$ terms are
the standard ones, unchanged). Such a $\nu$ lies in $\Ldual$: at $N$ by $L_N^\vee=L_{+,N}^\vee
\oplus L_{-,N}^\vee$, which is Lemma~\ref{lem:LN}, and elsewhere because $\nu_p\in L_{-,p}\subset
L_p$. So the condition $x+\nu\in\Ldual$ reads $x\in\Ldual\cap\Q\lambda$ and does not depend on
$\nu$. The $x=0$ term is $c_{[\nu]}(0)$, with $[\nu]$ a nonzero isotropic coset of $\Ldual/L$; the
$2N-2$ values agree by Lemma~\ref{lem:iso}. This gives the display, with the prefactor
$-\lvert\mathrm{CM}(d)\rvert/4$ of Theorem~B and $-\tfrac14(2N-2)(-\log N/(N-1))=\tfrac12\log N$;
the restored terms are \emph{added}, as \textup{Observation~obs:mult of the preprint} requires.

For $\kappa_\nu(0)$ itself take $\eta=\nu$ and $m=0$: only $x=0$ contributes, which needs
$\mu_+\in\Lp$ since $\nu_+=0$; then $\mu_-=\mu-\mu_+\in L\cap U=\Lm$, so $\mu=0$ in $L/(\Lp+\Lm)$
is the only class, and for it $\eta_-+\mu_-\equiv\nu$ modulo $\Lm$. Hence $\kappa_\nu(0)=
\kappa^-_\nu(0)$. (For $\eta=0$ the same count gives $\mu=0$ alone, which is \cite[Lemma~20]{Yang}'s
$\kappa_0(0)=\kappa^-_0(0)$, unchanged.)
\end{proof}

The same accounting is visible in Schofer's own reduction: the $(n,2)$ theorem is proved by
applying the $(0,2)$ theorem to $\langle F,\theta_+\rangle$, whose constant term at the coset
$\eta_-+\mu_-$ is $\sum_{x\in\eta_++\mu_++\Lp}c_\eta(-Q(x))$ \cite[Cor.~3.2]{Schofer}, and the
$m=0$ part of the $(0,2)$ proof is $2\sum_{\mu\neq0}c_\mu(0)b_\mu(0,t)$ with \emph{that}
constant term, before Lemma~2.21 sets it to zero.

\begin{lemma}[when (H) holds]\label{lem:H}
Let $F_f$ be built from a scalar form $f$ on $\Gamma_0(M)$ with poles only at the cusps $\infty$
and $0$, as in \cite[(6)]{GY}. By \cite[Lemma~24]{GY}, in the normalisation of \textup{Section~sec:data of the preprint}
(where it is checked numerically on every form used), the principal part of $F_f$ at a nonzero
coset $\eta$ is that of $f$ at the cusp $0$, $c_\eta(-n/M)$ being the coefficient of $q^{-n/M}$ in
$f|S$, and depends on $\eta$ only through $Q(\eta)\bmod\Z$; at $\eta=0$ the principal part of $f$
at $\infty$ is added. Hence:
\begin{enumerate}
\item Let $\Lp=\Z\lambda_0$, with $\lambda_0=\lambda$ or $\lambda/2$ according as $d\equiv1$ or
      $0\bmod4$ \cite[proof of Lemma~18]{Yang}, so $Q(\lambda_0)=\lvert d\rvert$ or
      $\lvert d\rvert/4$, and $\Ldual\cap\Q\lambda=\Z\lambda_0/g$ with $g$ the content of
      $\langle\lambda_0,\cdot\rangle$ on $L$: $g$ divides $2Q(\lambda_0)$, is prime to odd $N$ and
      to the primes $p\mid D$, $p\nmid d$, is divisible exactly by $2$ when $d\equiv1\bmod4$ or
      $2\mid D$, and exactly by $p$ for each odd $p$ dividing both $D$ and $d$ (there $p\mid c_1$
      and $c_2$ is a unit in Yang's basis). If the pole of $f$ at the cusp $0$ has order below
      $M\cdot\min_{j\ge1}j^2Q(\lambda_0)/g^2=MQ(\lambda_0)/g^2$ in $q^{1/M}$, then (H) holds at
      $d$. The threshold can be small: on $X_0^{15}(2)$ at $d=-15$, $g=30$ and it is $1/60$ in
      powers of $q$, i.e.\ a single power of the local parameter $q^{1/60}$.
\item In particular (H) holds at every $d$ when $f$ is holomorphic at the cusp $0$. A pole of
      order below $M$ (one power of $q$) removes only the terms with $x\in L$; those with
      $x\notin L$ sit at the exponents $Mj^2Q(\lambda_0)/g^2$, $g\nmid j$, which lie below $M$ as
      soon as $g$ is large.
\item For $x\notin L$ the coefficient $c_{[x+\nu]}(-Q(x))$ equals $c_{[x]}(-Q(x))$, both cosets
      being nonzero (the $N$-part of $x+\nu$ is not in $L_N$, the sum there being direct) with
      the same norm; for $x\in L$ it is the cusp-$0$ part $c_0(-Q(x))-c_\infty(-Q(x))$ of the
      coefficient at the zero coset.
\end{enumerate}
\end{lemma}

\begin{remark}[the $x\neq0$ terms under Theorem~B's own hypothesis]\label{rem:xsum}
Theorem~B is applied at $\tau_d\notin\mathrm{supp}\,\mathrm{div}\,\psi_F$. By Borcherds' theorem
in the form \cite[Thm.~A(2)]{Yang}, the order of $\Psi_F$ at $\tau_d$ is $\sum_{0<r,\,r\lambda\in
\Ldual}c_{r\lambda}(-r^2\langle\lambda,\lambda\rangle/2)=\sum_{x\in\Ldual\cap\Q\lambda,\,x>0}
c_{[x]}(-Q(x))$, and $c_{[-x]}=c_{[x]}$; so $\sum_{x\neq0}c_{[x]}(-Q(x))=0$ at an evaluated point,
and by Lemma~\ref{lem:H}(3)
\[
  \sum_{x\neq0}c_{[x+\nu]}(-Q(x))
   =\sum_{x\neq0}c_{[x]}(-Q(x))-\sum_{x\in L\cap\Q\lambda,\,x\neq0}c_\infty(-Q(x))
   =-2\sum_{k\ge1}c_\infty\bigl(-k^2Q(\lambda_0)\bigr),
\]
with $c_\infty(-n)$ the coefficient of $q^{-n}$ in $f$ at $\infty$. So the full correction per CM
point is
\[
  \log N\Bigl(\tfrac12c_\eta(0)-\sum_{k\ge1}c_\infty\bigl(-k^2Q(\lambda_0)\bigr)\Bigr):
\]
the $x\neq0$ terms are nonzero exactly when $f$ has a pole at $\infty$ of order $k^2Q(\lambda_0)$
whose contribution to the order at $\tau_d$ is cancelled by the cusp-$0$ side. Such a point is a
legitimate point of evaluation --- Theorem~B's hypothesis is the vanishing of the order, not of
each term --- with a correction different from $\tfrac12c_\eta(0)$ and dependent on $d$. Under
Lemma~\ref{lem:H}(1), in particular when $f$ is holomorphic at the cusp $0$, no such cancellation
is possible; otherwise the sum has to be evaluated at each point, which costs a few coefficient
lookups, and the accompanying code now does so (\texttt{M0InfinityPoleSum}).

On the data: the input forms of every even-$D$ base are holomorphic at the cusp $0$, so (H) holds
there outright. On the six odd-$D$ bases with prime $N>1$ in our model set
($X_0^{15}(2)$, $X_0^{21}(2)$, $X_0^{35}(2)$, $X_0^{39}(2)$, $X_0^{51}(2)$, $X_0^{57}(2)$) the
pole at the cusp $0$ has order up to $\tfrac54M$ (on $X_0^{39}(2)$), so (H) can fail; there the
order at $\tau_d$ and the $x\neq0$ sum were computed from the two cusp expansions for every form
and every fundamental $d$ at which the term fires, and the $x\neq0$ sum was zero at every point
off the divisor. This is a structural fact of these forms rather than a tested cancellation --- the
principal part of each at the cusp $0$ is a single monomial whose exponent is never of the form
$Mj^2Q(\lambda_0)/g^2$ --- and on $X_0^{39}(2)$ and $X_0^{51}(2)$ the forms as returned by the
search have non-integral principal parts at $\infty$, outside Theorem~B's hypothesis
$c_\eta(-m)\in\Z$, so those two bases are not evidence either way. The agreement of
\textup{Observation~obs:mult of the preprint} could therefore not have detected the $x\neq0$ terms; it is
Proposition~\ref{prop:mult} that accounts for them.
\end{remark}

The same argument settles the two cases the hypotheses exclude.

\begin{proposition}[composite squarefree level]\label{prop:composite}
Let $N$ be squarefree, $d$ fundamental, and for a nonzero isotropic $\nu\in\Lm^\vee/\Lm$ let
$S(\nu)=\{p\mid N:\nu_p\neq0\}$; then $S(\nu)$ is nonempty and consists of primes not dividing
$d$, and
\[
  \kappa^-_\nu(0)=\begin{cases}-\log p/(p-1)&\text{if }S(\nu)=\{p\},\\ 0&\text{if }\lvert S(\nu)
  \rvert\ge2.\end{cases}
\]
Consequently, under (H) at every prime, the restored terms of Theorem~B contribute
$\sum_{p\mid N,\,p\nmid d}\tfrac12c_{\eta_p}(0)\log p$ per CM point, where $c_{\eta_p}(0)$ is the
common constant term of the $2p-2$ cosets supported exactly at $\{p\}$. Of the $2^{\omega(N)}-1$
functionals of \textup{Section~sec:determined of the preprint}, the multiplier at the prime $p$ is the one attached to the
support class $\{p\}$; the classes supported at two or more primes do not enter.
\end{proposition}
\begin{proof}
At $p\mid N$ with $p\nmid d$ the local plane is the $p$-scaled split plane (Lemma~\ref{lem:LN},
prime by prime; Eichler's condition makes $p$ split), and at every other prime --- $p\nmid N$, or
$p\mid N$ with $p\mid d$ --- every nonzero local coset is anisotropic (Lemma~\ref{lem:Wp} and
Remark~\ref{rem:ramified}). So $S(\nu)$ is as stated, and for $\mathrm{Re}\,s$ large
\[
  E_0(\tau,s;\varphi_\nu,+1)=\Phi(s)\prod_{p\in S(\nu)}\frac{1-p^{-s}}{1+(p-2)p^{-s}}
\]
by Lemma~\ref{lem:WN} at each $p\in S(\nu)$, with $\Phi(0)=-1$ as in the proof of
Proposition~\ref{prop:kappa0} (the poles of $W_{0,p}(s,\varphi_0)$ at the primes of $N$ are
cancelled by the missing Euler factors of $L^S(s,\chi_d)$). The product has a zero of order
$\lvert S(\nu)\rvert$ at $s=0$, whose first-order coefficient is $\log p/(p-1)$ when
$S(\nu)=\{p\}$ and $0$ otherwise. The count $\kappa_\nu(0)=\kappa^-_\nu(0)$ is as in
Proposition~\ref{prop:mult}: with $\eta=\nu$ and $m=0$, $\mu_+\in\Lp$ forces $\mu=0$. The cosets
with support exactly $\{p\}$ are the nonzero isotropic vectors of the plane at $p$ with zero
components elsewhere, one orbit of its orthogonal group extended by the identity, so they share
$c_{\eta_p}(0)$ (Lemma~\ref{lem:iso}); there are $2p-2$ of them, and
$-\tfrac14(2p-2)(-\log p/(p-1))=\tfrac12\log p$. The $x\neq0$ terms are handled prime by prime
exactly as in Lemma~\ref{lem:H} and Remark~\ref{rem:xsum}.
\end{proof}

\begin{remark}[a prime of $N$ dividing $d$]\label{rem:ramified}
If $p\mid N$ and $p\mid d$, nothing is dropped at $p$: $L_{-,p}$ is anisotropic, so no nonzero
isotropic coset is supported there and $W_{0,p}(s,\varphi_{\nu,p})=0$ for every $\nu_p\neq0$.
For odd $p$ this is \cite[Lemma~18(1)]{Yang} with $p\mid(DN,d)$, Gram
$\mathrm{diag}(\epsilon_1,\epsilon_2)$ with $p\,\|\,\epsilon_2$, as in the proof of
Lemma~\ref{lem:Wp}. For $p=2$, a case Yang omits, write $\lambda_0=\lambda/2=\begin{pmatrix}c_1&
c_2\\2c_3&-c_1\end{pmatrix}$ in $\mathcal O\otimes\Z_2=\begin{pmatrix}\Z_2&\Z_2\\2\Z_2&\Z_2
\end{pmatrix}$, $c_1^2+2c_2c_3=d/4\equiv2,3\bmod4$, and $L_{-,2}=\{\begin{pmatrix}a&b\\2c&-a
\end{pmatrix}:ac_1+bc_3+cc_2=0\}$ with $Q=-(a^2+2bc)$ (the reduced norm; only its integrality
is used below). If $d/4\equiv3\bmod4$ then $c_1,c_2,c_3$ are
odd; eliminating $a$ gives a form $2(uy^2+2v\,yz+u'z^2)$ with $u,u'$ units and $v$ even, hence
$2\,\mathrm{diag}(\epsilon_1,\epsilon_2)$ with $\epsilon_1\epsilon_2=-(d/4)/c_1^2\equiv1\bmod4$,
and the three nonzero cosets of $(\Z/2)^2$ have $Q\equiv\epsilon_1/4,\epsilon_2/4,\epsilon_1/2$,
none integral. If $d/4\equiv2\bmod4$ then $c_1$ is even and $c_2,c_3$ odd; eliminating $b$ gives
$2(a^2-2(c_1/c_3)ac-2(c_2/c_3)c^2)$, hence $2\,\mathrm{diag}(1,2\epsilon')$ with
$\epsilon'$ a unit, whose discriminant group $\Z/2\times\Z/4$ has $Q(a/2,b/4)=a^2/4+\epsilon'b^2/8$,
never integral on a nonzero element. This is the firing rule $\sum_{p\mid N/(N,d)}\log p$: the
primes of $N$ dividing $d$ contribute nothing, for the reason the recipe assumed at all of them.
\end{remark}


\subsection*{Non-fundamental discriminants}

Everything above assumes $d$ fundamental. Write $d=f^2d_0$ with $d_0$ fundamental and $f>1$ the
conductor, $R_f=\Z+f\mathcal O_{d_0}$, and let $\lambda=\varphi(\sqrt d)$ for an optimal embedding
$\varphi$ of $R_f$, so that $\langle\lambda,\lambda\rangle=2\lvert d\rvert$ as before. At a level prime $N$ dividing $f$ the lattice is different, and so is the bookkeeping;
what is proved below is the general form of the restored terms (Proposition~\ref{prop:fibre})
and, at an odd conductor prime, the local factors that enter it (Lemma~\ref{lem:Wcond}); at $p=2$
the factors are computed on $X_0^{15}(2)$ and agree with the same formulas.

\begin{lemma}[the plane's isotropy is a property of the field]\label{lem:split}
For every prime $p$, the quadratic space $\Lm\otimes\Q_p=\lambda^\perp\subset B_0\otimes\Q_p$ is
isotropic if and only if $p$ splits in $k=\Q(\sqrt{d_0})$. In particular this does not depend on
the conductor.
\end{lemma}
\begin{proof}
Write $B=k\oplus kj$ with $j$ anticommuting with $k$ and $j^2=\beta\in\Q^\times$, and $B_0$ for the
trace-zero part. Then $\lambda\in k_0=\Q\sqrt{d_0}$ and $\lambda^\perp=kj$: for $x\in k$,
$\tr(\lambda\overline{xj})=\tr(\lambda\,\bar j\,\bar x)=-\tr(\lambda\bar x j)=0$ since $\lambda\bar x
\in k$ and $\tr(yj)=0$ for $y\in k$. The quadratic form is $Q(xj)=\nrd(xj)=-\beta\,N_{k/\Q}(x)$.
So $\lambda^\perp\otimes\Q_p\cong(k_p,-\beta N)$, which represents $0$ nontrivially exactly when
$N_{k_p/\Q_p}$ does, i.e.\ when $k_p$ is not a field.
\end{proof}

\begin{lemma}[pole or no pole at the zero coset]\label{lem:pole}
In the counting normalisation of \textup{Lemma~lem:m0loc of the preprint}, $W_{0,p}(s,\varphi_0)=
(1-X)\sum_{k\ge0}\alpha_kX^k$ with $X=p^{-s}$, $\alpha_0=1$ and $\alpha_k=p^{-k}\#\{x\in
L_{-,p}/p^kL_{-,p}:\ Q(x)\equiv0\bmod p^k\}$. If $L_{-,p}\otimes\Q_p$ is anisotropic, the $\alpha_k$
are bounded and $W_{0,p}(s,\varphi_0)$ is holomorphic at $s=0$. If it is isotropic, $\alpha_k\ge
ck$ for some $c>0$ and $W_{0,p}(s,\varphi_0)$ has a pole at $s=0$.
\end{lemma}
\begin{proof}
Anisotropic: $Q$ is a nonzero anisotropic form on $\Z_p^2$, so there is $c_0$ with
$\ord_p Q(x)\le c_0+2\min_i\ord_p(x_i)$ for all $x\ne0$ (on the compact set of primitive $x$,
$\ord_pQ(x)$ is bounded, and $Q(p^ax)=p^{2a}Q(x)$). Hence $Q(x)\equiv0\bmod p^k$ forces every
coordinate of $x$ to vanish modulo $p^{\lceil(k-c_0)/2\rceil}$, which leaves at most $p^{k+c_0}$
classes modulo $p^k$, and $\alpha_k\le p^{c_0}$. Then $\lvert(1-X)\sum\alpha_kX^k\rvert\le p^{c_0}$
for $0<X<1$, so there is no pole. Isotropic: choose a primitive isotropic $e\in L_{-,p}$ and
complete it to a basis $e,f'$, so $Q(xe+yf')=uxy+wy^2$ with $u=\langle e,f'\rangle\ne0$. For each
$0\le j\le k-\ord_pu-1$, the classes $y$ with $\ord_py=j$ number $(p-1)p^{k-j-1}$, and for each of
them $Q\equiv0\bmod p^k$ is a congruence on $x$ modulo $p^{k-j-\ord_pu}$ with $p^{j+\ord_pu}$
solutions modulo $p^k$; summing, $\alpha_k\ge(k-\ord_pu)(p-1)/p$, which grows linearly, and the
pole follows since $(1-X)\sum\alpha_kX^k\ge(1-X)\sum_{k\le K}ckX^k$ for every $K$.
\end{proof}

\begin{lemma}[the lattice at a conductor prime]\label{lem:conductor}
Let $p$ be an odd level prime, $p\,\|\,N$, $p\nmid d_0$, $p^k\,\|\,f$ with $k\ge1$, and
$\lambda$ the image of $\sqrt d=f\sqrt{d_0}$ under an optimal embedding of $R_f$. In the model
$\mathcal O\otimes\Z_p=\begin{pmatrix}\Z_p&\Z_p\\p\Z_p&\Z_p\end{pmatrix}$, $L_p=\{\begin{pmatrix}
a&b\\pc&-a\end{pmatrix}\}$ with $Q=\nrd$, so $Q(a,b,c)=-a^2-pbc$:
\begin{enumerate}
\item $\lambda=\begin{pmatrix}pa'&b\\p^2c'&-pa'\end{pmatrix}$ with $b$ a unit and
$a'^2+bc'=-\lvert d\rvert/p^2$;
\item $L_{-,p}$ has the basis $e_1=(1,0,-2a'/b)$, $e_2=(0,1,-pc'/b)$ and, after the integral change
of basis $e_1'=e_1$, $e_2'=e_2-(pa'/b)e_1$, the Gram form
$Q=-\alpha^2-p^{2k}\lvert d_0\rvert v^2\beta^2$ with $v$ a unit: a unit line orthogonal to a
$p^{2k}$-scaled line;
\item $[L_p:\Z_p\lambda\oplus L_{-,p}]=p^{2k-1}$;
\item $L_{-,p}^\vee/L_{-,p}\cong\Z/p^{2k}$, generated by $e_2'/p^{2k}$; its cosets of integral norm are
the $p^k-1$ multiples of $e_2'/p^k$, and of these exactly the $p-1$ multiples of $e_2'/p$ lie in
$L_p^\vee$. In particular the direct-sum decomposition of Lemma~\ref{lem:LN} fails, and a coset of
$L_{-,p}^\vee/L_{-,p}$ enters the decomposition of some $\varphi_\eta$, $\eta\in\Ldual/L$, only if it
is one of those $p-1$; the others can pair only with $x\notin\Ldual$, $x+\nu\in\Ldual$
(Proposition~\ref{prop:fibre}).
\end{enumerate}
For $p=2$ the same computation on $X_0^{15}(2)$ gives, for $f=2,4,8$ and both splitting types,
index $2^{2k-1}$, elementary divisors of $2$-adic valuations $(1,2k-1)$, $2^k-1$ integral cosets of
which exactly one lies in $L_2^\vee$; the pole at the zero coset is present exactly when $2$ splits.
\end{lemma}
\begin{proof}
Write $\lambda=(a,b,c)$. Then $Q(\lambda)=-a^2-pbc=\lvert d\rvert\cdot(\text{unit})^2$ is divisible by
$p^{2k}$ with $2k\ge2$, so $p\mid a$. Optimality of the embedding of $R_f$ says $\lambda/p\notin
\mathcal O\otimes\Z_p$, i.e.\ not both $p\mid a$ and $p\mid b$, so $b$ is a unit; then
$p^2\mid a^2$ forces $p\mid c$, which is (1). For $x=(\alpha,\beta,\gamma)\in L_p$,
$\langle x,\lambda\rangle=-\tr(x\lambda)=-(2\alpha pa'+\beta p^2c'+\gamma pb)$, so $x\perp\lambda$
iff $\gamma=-(2a'\alpha+pc'\beta)/b$, which gives the basis in (2); substituting,
$Q|_{L_{-,p}}=-\alpha^2+(2pa'/b)\alpha\beta+(p^2c'/b)\beta^2$, and completing the square gives
$-\alpha'^2+(p^2/b^2)(a'^2+bc')\beta^2=-\alpha'^2-p^{2k}\lvert d_0\rvert v^2\beta^2$ by (1). For
(3) the determinant of the matrix with rows $\lambda,e_1,e_2$ is $(2p/b)(a'^2+bc')$, of valuation
$2k-1$. For (4), $\langle e_2',e_2'\rangle=-2p^{2k}\lvert d_0\rvert v^2$ and $\langle e_1',e_1'\rangle
=-2$ are the only nonzero pairings, so the dual is spanned by $e_1'/2$ and $e_2'/(2p^{2k}\lvert d_0
\rvert v^2)$, i.e.\ by $e_2'/p^{2k}$ up to units; $Q(je_2'/p^{2k})=-\lvert d_0\rvert v^2j^2/p^{2k}$ is
integral iff $p^k\mid j$. Finally $\langle e_2'/p^m,x\rangle=-(1/p^{m-1})(\gamma+(pc'/b)\beta)+
\dots$ is integral on $L_p$ iff $m\le1$, which leaves the multiples of $e_2'/p$.
\end{proof}

So at a level prime dividing the conductor the picture splits by the field: if $p$ splits, the
plane is isotropic and the zero coset has its pole; if $p$ is inert, the plane is anisotropic and
there is no pole --- a case only non-fundamental $d$ can produce, since for fundamental $d$ an inert
$p\,\|\,N$ admits no optimal embedding (Eichler's condition $\nu_p=1+(\tfrac dp)=0$), whereas for
$p\mid f$ the local embedding number is $2$ regardless. In neither case does the argument of
Proposition~\ref{prop:mult} apply as it stands, because it identifies a coset of $\Lm^\vee/\Lm$ with
a coset of $\Ldual/L$ through $L_N=L_{+,N}\oplus L_{-,N}$, and Lemma~\ref{lem:conductor}(3) says the
sum is not direct. What survives is the bookkeeping before that identification.

\begin{proposition}[the restored terms without the direct-sum hypothesis]\label{prop:fibre}
Let $N$ be prime and $d<0$ any discriminant with $\mathrm{CM}(d)\neq\varnothing$. The terms of
\cite[Thm.~B]{GY} in which $\kappa^-$ is evaluated at $0$ at a nonzero coset total, per CM point,
\[
  -\frac14\sum_{\substack{\nu\in\Lm^\vee/\Lm\\ \nu\neq0}}\kappa^-_\nu(0)
     \sum_{\substack{x\in\Lp^\vee\\ x+\nu\in\Ldual}}c_{[x+\nu]}\bigl(-Q(x)\bigr),
\]
where $\kappa^-_\nu(0)=0$ unless $Q(\nu)\in\Z$. If moreover $\nu$ is supported at the single prime
$p$ (i.e.\ $\nu_q=0$ for $q\neq p$), then with $X=p^{-s}$ and $W_{0,p}(s,\varphi_\nu)$ the local
factor at $p$ in any fixed normalisation,
\[
  \kappa^-_\nu(0)=\log p\cdot\frac{d}{dX}\Bigl[\frac{W_{0,p}(s,\varphi_\nu)}{W_{0,p}(s,\varphi_0)}
  \Bigr]_{X=1},
\]
the ratio being a rational function of $X$ that vanishes at $X=1$.
\end{proposition}
\begin{proof}
The display is the first display in the proof of Proposition~\ref{prop:mult}, which is obtained
from Theorem~B by regrouping alone, before $L_N=L_{+,N}\oplus L_{-,N}$ is used; the prefactor is
Theorem~B's $-\tfrac14$ per point. For $Q(\nu)\notin\Z$ some $p$ has $Q(\nu_p)\notin\Z_p$, and then
$W_{0,p}(s,\varphi_\nu)\equiv0$ (\cite[Thm.~4.3(i), 4.4(i)]{KY}: every $x\in\nu_p+L_{-,p}$ has
$Q(x)\equiv Q(\nu_p)\notin\Z_p$), so $E_0(\tau,s;\varphi_\nu)\equiv0$. For $\nu$ supported at $p$,
exactly as in the proof of Proposition~\ref{prop:kappa0}, $E_0(\tau,s;\varphi_\nu)=\Phi(s)\,
W_{0,p}(s,\varphi_\nu)/W_{0,p}(s,\varphi_0)$ with $\Phi(0)=-1$; the series $E(\tau,s;\varphi_\nu)$ is
incoherent, hence vanishes at $s=0$, so the ratio vanishes at $X=1$, and $dX/ds=-X\log p$ gives the
formula. (For $\nu$ supported at two or more primes the product of the vanishing ratios has a
double zero and $\kappa^-_\nu(0)=0$, as in Proposition~\ref{prop:composite}.)
\end{proof}

At a conductor prime the inner sum is genuinely a sum over a fibre. Two kinds of pairs $(x,\nu)$
appear that do not exist for fundamental $d$: a coset $\nu\in\Ldual$ pairs with $x\notin L$ such
that $x+\nu\in L$ (so $[x+\nu]=0$ and the coefficient is $c_0(-Q(x))$, which includes the pole of $f$
at $\infty$), and a coset $\nu\notin\Ldual$ pairs with $x\notin\Ldual$ such that $x+\nu\in\Ldual$.
The vectors $x$ that occur are the vectors $\lambda_0/p^j$ of the CM line, of norm
$Q(\lambda_0)/p^{2j}$: the primitive CM vectors of the discriminants $d/p^{2j}$. So the correction
at a conductor prime sees the poles of $f$ at the CM points of the \emph{smaller} discriminants
$d/p^2,d/p^4,\dots$, which no $d$-independent multiplier can.

\paragraph{Computation on $X_0^{15}(2)$.} Everything in Proposition~\ref{prop:fibre} is computable
from the lattice: the integral cosets of $\Lm^\vee/\Lm$ supported at $2$, the fibre over each, the
coefficients $c_{[x+\nu]}(-Q(x))$ (the cusp-$0$ principal part of $f$ at the nonzero cosets by
Lemma~\ref{lem:H}, $c_\eta(0)=2\,\mathrm{mult}(f)$ at a nonzero isotropic coset, and the
cusp-$\infty$ principal part added at $\eta=0$), and the local factors $W_{0,2}(s,\varphi_\nu)=
(1-X)\sum_k\alpha_kX^k$ by counting on $\Lm=L\cap\lambda^\perp$ as in Lemma~\ref{lem:pole},
reconstructed as rational functions from $\alpha_0,\dots,\alpha_{10}$ (\texttt{level-p2/fibresum.m}
on the campaign branch; every reconstruction has at least two spare terms). With $\lambda_0=
\lambda/2$ ($Q(\lambda_0)=\lvert d\rvert/4$) and $\nu_0$ the integral coset lying in $\Ldual$:
\begin{center}\small
\begin{tabular}{lllll}
$d$ & $f$, type & $W_{0,2}(s,\varphi_0)$ & integral cosets $\nu$: $W_{0,2}(s,\varphi_\nu)$, $\kappa^-_\nu(0)/\log2$ & fibre over $\nu$ ($Q(x)$) \\ \hline
$-15,-7$ & $1$, split & $1/(1-X)$ & two, in $\Ldual$: $1$, $-1$ & $x=0$ \\
$-60,-28$ & $2$, split & $(1-X+X^2)/(1-X)$ & $\nu_0$: $1+X$, $-2$ & $x=0$ \\
$-240$ & $4$, split & $(1-X+X^2-X^3+2X^4)/(1-X)$ & $\nu_0$: $1+X^2+2X^3$, $-2$ & $x=0$;\ $x=\pm\lambda_0/2$, $x+\nu_0\in L$ ($15$) \\
 & & & $\pm\nu_1\notin\Ldual$: $1+X$, $-1$ & none below the poles \\
$-12$ & $2$, inert & $(1+X+X^2)/(1+X)$ & $\nu_0$: $1-X$, $-\tfrac23$ & $x=0$;\ $x=\pm\lambda_0/2$, $[x+\nu_0]\ne0$ ($\tfrac34$) \\
$-48$ & $4$, inert & $(1+X+X^2+X^3+2X^4)/(1+X)$ & $\nu_0$: $(1-X)(1+X+2X^2)$, $-\tfrac43$ & $x=0$;\ $x=\pm\lambda_0/2$, $x+\nu_0\in L$ ($3$) \\
 & & & $\pm\nu_1\notin\Ldual$: $1-X$, $-\tfrac13$ & $x=\pm\lambda_0/4$, $[x+\nu_1]\neq0$ ($\tfrac34$) \\
\end{tabular}
\end{center}
(The fibres are listed up to the largest pole of the nine forms of the model set, $q^{-30}$ at
$\infty$ and $q^{-3/4}$ at the cusp $0$.) Writing $m=\tfrac12c_\eta(0)$ for the fundamental
multiplier, $a=c_\infty(-\lvert d\rvert/16)$ for the coefficient of $f$ at the pole exponent
$Q(\lambda_0/2)$ and $b$ for the coefficient of $q^{-3/4}$ in $f|S$, the proposition gives per point
\[
  d=-240:\quad (m+a)\log2,\qquad\qquad d=-48:\quad \bigl(\tfrac23(m+a)+\tfrac13 b\bigr)\log2 ,
\]
while at $d=-60,-28$ it gives $m\log2$ (one coset at twice the weight of the two fundamental ones)
and at $d=-15,-7$ the $m\log2$ of Proposition~\ref{prop:mult}. These predictions were compared, for
all nine forms, with the values forced by Table~45: each form is a function on the genus-$0$ curve
$X_0^{15}(2)/W$ whose divisor is known (Lemma~25 of \cite{GY}), so its value at every CM point is
$C\prod_i\lvert s(d)-s(d_i)\rvert^{m_i}$ with $s$ Guo--Yang's Hauptmodul; the constant $C$ was read
off at $d=-7$ and the formula checked at six to eight further fundamental and conductor-$2$ points per
form before being used at $-240$ and $-48$ (\texttt{level-p2/evidence.log}). The fibre sum
reproduces all $18$ values exactly. The rule ``add $m\log p$ when $p$ splits, nothing when $p$ is
inert'', which the routine applied until now, is right at $-240$ only for the five forms with
$a=0$ (those whose divisor misses $\tau_{-60}$) and at $-48$ only for the three forms with
$m+a=0$ and $b=0$ --- among them the form on which the rule had been tested --- and Yang's conductor
term for the pairs with $Q(x)=m$, which the routine also carried, gives the irrational $2^{4/3}$ at
$-48$ for the two forms with $b\neq0$. The vanishing at $-48$ for the tested form is thus the
identity $c_\eta(0)+2c_\infty(-3)=0$ for that form, not the absence of a term.

The local factors in the table are instances of a closed form.

\begin{lemma}[the $m=0$ local factors at a conductor prime]\label{lem:Wcond}
Let $p$ be odd, $p\,\|\,N$, $p^k\,\|\,f$ with $k\ge1$, $\epsilon=(\tfrac{d_0}{p})=\pm1$, and
$L_{-,p}=\langle-1\rangle\perp\langle-p^{2k}c\rangle$ with $c$ a unit as in
Lemma~\ref{lem:conductor}. In the counting normalisation $W_{0,p}(s,\varphi_\nu)=(1-X)\sum_{j\ge0}
\alpha_j(\nu)X^j$, $X=p^{-s}$, $\alpha_j(\nu)=p^{-j}\#\{x\in(\nu+L_{-,p})/p^jL_{-,p}:Q(x)\equiv0
\bmod p^j\}$:
\begin{enumerate}
\item The cosets of integral norm are $\nu_r=(0,r/p^k)$, $r\bmod p^k$. For $r\neq0$ let
      $\rho=\ord_p r<k$ (so $\nu_r\in L_p^\vee$ iff $\rho=k-1$). Then $\alpha_j(\nu_r)=
      p^{\lfloor j/2\rfloor}$ for $j\le2\rho$ and $\alpha_j(\nu_r)=(1+\epsilon)p^\rho$ for
      $j>2\rho$, i.e.
      \[
        W_{0,p}(s,\varphi_{\nu_r})=(1-X)\sum_{j=0}^{2\rho}p^{\lfloor j/2\rfloor}X^j
                                   +(1+\epsilon)\,p^\rho X^{2\rho+1}.
      \]
\item For the zero coset, $\alpha_j(0)=p^{\lfloor j/2\rfloor}$ for $j\le2k$ and
      $\alpha_j(0)=(1+\epsilon)(p-1)p^{k-1}\lceil(j-2k)/2\rceil+p^{\,k-(j\bmod2)}$ for $j>2k$, i.e.
      \[
        W_{0,p}(s,\varphi_0)=(1-X)\Bigl[\sum_{j=0}^{2k}p^{\lfloor j/2\rfloor}X^j
          +\frac{X^{2k+1}(p^{k-1}+p^kX)}{1-X^2}\Bigr]
          +\frac{(1+\epsilon)(p-1)p^{k-1}X^{2k+1}}{1-X^2}.
      \]
\item Consequently $\kappa^-_{\nu_r}(0)=-\dfrac{2p^{\rho-k+1}}{p-1}\log p$ if $p$ splits, and
      \[
        \kappa^-_{\nu_r}(0)=-\frac{2}{p^{k-1}(p+1)}\Bigl(p^\rho+\frac{2(p^\rho-1)}{p-1}\Bigr)\log p
      \]
      if $p$ is inert; all other nonzero cosets have $\kappa^-_\nu(0)=0$.
\end{enumerate}
\end{lemma}
\begin{proof}
With $Q(u,t)=-(u^2+c\,p^{2k}t^2)$ the dual is generated by $e_2/p^{2k}$, and $Q(0,r/p^{2k})=
-cr^2/p^{2k}$ is integral iff $p^k\mid r$, which is (1)'s list and Lemma~\ref{lem:conductor}(4).
A vector of $\nu_r+L_{-,p}$ is $(u,w)$ with $w=r+p^kt\in r+p^k\Z_p$, and for $r\ne0$ every such $w$
has $\ord_pw=\rho$. The congruence $u^2+cw^2\equiv0\bmod p^j$ reads $u^2\equiv0$ when $j\le2\rho$,
with $p^{\lfloor j/2\rfloor}$ solutions $u\bmod p^j$ for each of the $p^j$ values of $t$; when
$j>2\rho$ it forces $u=p^\rho u'$ with $u'^2\equiv-c(w/p^\rho)^2\bmod p^{j-2\rho}$, which has $2$
solutions $u'\bmod p^{j-2\rho}$ if $-c$ is a square mod $p$ --- i.e.\ $\epsilon=1$, since
$-c=d_0v^2$ --- and none otherwise, giving $(1+\epsilon)p^{\rho}$ values of $u\bmod p^j$ for each
$t$. This is (1). For the zero coset, $w=p^kt$; if $t\equiv0\bmod p^j$ or $2k+2\ord_pt\ge j$ the
congruence is $u^2\equiv0$ with $p^{\lfloor j/2\rfloor}$ solutions, and the $t$ with
$2k+2v<j$, $v=\ord_pt$, number $(p-1)p^{j-v-1}$ and each admits $(1+\epsilon)p^{k+v}$ values of
$u$ by the same Hensel count. Summing, for $j>2k$,
\[
  p^j\alpha_j(0)=(1+\epsilon)(p-1)p^{j+k-1}\Bigl\lceil\frac{j-2k}2\Bigr\rceil
   +p^{\lfloor j/2\rfloor}\Bigl(1+\sum_{v\ge\lceil(j-2k)/2\rceil}^{j-1}(p-1)p^{j-1-v}\Bigr)
   =(1+\epsilon)(p-1)p^{j+k-1}\Bigl\lceil\frac{j-2k}2\Bigr\rceil+p^{\lfloor j/2\rfloor+j-\lceil(j-2k)/2\rceil},
\]
and $\lfloor j/2\rfloor-\lceil (j-2k)/2\rceil=k-(j\bmod2)$; for $j\le2k$ only $u^2\equiv0$ remains.
The generating functions follow from $\sum_{n\ge1}\lceil n/2\rceil X^n=X/((1-X)(1-X^2))$. For (3)
apply Proposition~\ref{prop:fibre}. If $\epsilon=-1$, $W_{0,p}(s,\varphi_0)$ is finite at $X=1$ with
value $(p^{k-1}+p^k)/2$ (only the middle term survives the factor $1-X$) and $W_{0,p}(s,
\varphi_{\nu_r})=(1-X)\sum_{j\le2\rho}p^{\lfloor j/2\rfloor}X^j$ has derivative
$-\sum_{j\le2\rho}p^{\lfloor j/2\rfloor}=-(p^\rho+2(p^\rho-1)/(p-1))$ there. If $\epsilon=1$,
$(1-X)W_{0,p}(s,\varphi_0)\to(p-1)p^{k-1}$ at $X=1$ while $W_{0,p}(s,\varphi_{\nu_r})\to2p^\rho$,
so the ratio has derivative $-2p^\rho/((p-1)p^{k-1})$.
\end{proof}

The formulas were also checked by brute-force counting on the plane $L\cap\lambda^\perp$ of
$X_0^{10}(3)$ at $p=3$ --- $d=-72,-648$ (split, $k=1,2$) and $d=-387,-603$ (inert, $k=1$), every
$3$-primary integral coset, six terms each (\texttt{level-p2/wcond\_check.m}). At $p=2$ the same
formulas with $p=2$ reproduce every series computed on $X_0^{15}(2)$ --- the
six rows of the table, $k=1,2$, both splitting types --- although the $2$-adic plane is only
computed to be of the shape of Lemma~\ref{lem:conductor}; the factors $-2$, $-1$, $-\tfrac43$,
$-\tfrac13$, $-\tfrac23\log2$ are the values of (3) at $(k,\rho)=(1,0),(2,1),(2,0)$ in the split
case and $(2,1),(2,0),(1,0)$ in the inert case. With Lemma~\ref{lem:Wcond}, Proposition~\ref{prop:fibre}
is explicit at an odd conductor prime: the restored terms there are
\[
  -\frac14\sum_{r\neq0}\kappa^-_{\nu_r}(0)\sum_{\substack{x\in\Lp^\vee\\ x+\nu_r\in\Ldual}}
  c_{[x+\nu_r]}\bigl(-Q(x)\bigr),
\]
with the fibre over each $\nu_r$ read off the lattice as in the computation above. What remains for
$p=2$ is the lattice itself (the $2$-adic part of Lemma~\ref{lem:conductor}); the routine should
evaluate the sum as written rather than apply a rule, and the values in the table are what it must
reproduce.

\begin{remark}[the $m>0$ terms at non-fundamental $d$]
A separate point, settled by the same comparison: in the $m>0$ part of the formula the class
number and unit count that enter are those of the field $k$, not of the order $R_f$, together
with a factor $1/f$. The two agree exactly when $h(R_f)=h(k)$, which holds at conductor $2$ when
$2$ splits ($d=-28,-60$) and fails at conductor $4$, where $h(R_f)=2h(k)$ and the computed values
came out halved against Table~45 down to a half-integer exponent of $5$; with the field's class
number $d=-240$ is exact and nothing at a fundamental discriminant moves. In the invariants of
the order this factor reads $(w_{R_f}/h(R_f))\prod_{p\mid f}(1-\chi_{d_0}(p)/p)$.
\end{remark}

\begin{thebibliography}{9}
\bibitem{Schofer} J.~Schofer, \emph{Borcherds forms and generalizations of singular moduli},
J.\ reine angew.\ Math.\ \textbf{629} (2009), 1--36.

\bibitem{SchoferThesis} J.~Schofer, \emph{Borcherds forms and generalizations of singular moduli},
Ph.D.\ thesis, University of Maryland, College Park, 2005.

\bibitem{KY} S.~Kudla and T.~Yang, \emph{Eisenstein series for $SL(2)$},
Sci.\ China Math.\ \textbf{53} (2010), 2275--2316.

\bibitem{Yang} Y.~Yang, \emph{Special values of hypergeometric functions and periods of CM elliptic
curves}, arXiv:1503.07971.

\bibitem{GY} J.-W.~Guo and Y.~Yang, \emph{Equations of hyperelliptic Shimura curves},
Proc.\ Lond.\ Math.\ Soc.\ (2017), doi:10.1112/S0010437X16007739.

\bibitem{Borcherds} R.~E.~Borcherds, \emph{Automorphic forms with singularities on Grassmannians},
Invent.\ Math.\ \textbf{132} (1998), 491--562.
\end{thebibliography}

\end{document}
